\documentclass[indonesian,a4paper,oneside,11pt,titlepage]{hitec}
\usepackage[sfdefault]{FiraSans}
\usepackage{FiraMono}
\usepackage[utf8]{inputenc}
\usepackage{awesomebox}
\usepackage{listings}
\usepackage[indonesian]{babel}
\usepackage{blindtext}
\usepackage{cmap}
\usepackage{graphicx}
\usepackage[pdftitle={Panduan Pengguna dan Administrasi e-SPPO V2.0.0},pdfdisplaydoctitle]{hyperref}
\usepackage{nameref}
\usepackage{bookmark}

\title{Panduan Pengguna dan Administrasi\linebreak Sistem Penyelenggaraan Pemilihan OSIS Elektronik (e-SPPO)\linebreak Versi 2.0.0}
\author{Rizki Yandri}
\company{SMAN 1 Bati-Bati -- Politeknik Negeri Tanah Laut}
\date{Agustus 2025}

\begin{document}

\maketitle
\newpage
\setcounter{secnumdepth}{3}
\setcounter{tocdepth}{3}
\tableofcontents

% ===================================================================================
% BAB 1: PENDAHULUAN
% Berisi pengenalan umum mengenai metode pemilihan OSIS dan aplikasi e-SPPO.
% File ini akan di-include oleh main.tex
% ===================================================================================

\part{Pendahuluan}
\label{part:pendahuluan}

Dokumen ini merupakan panduan resmi untuk penggunaan dan administrasi \textbf{Sistem Penyelenggaraan Pemilihan OSIS Elektronik (e-SPPO)} versi 2.0.0. Panduan ini dirancang untuk menjadi referensi utama bagi seluruh pihak yang terlibat dalam pemilihan, termasuk panitia penyelenggara, pengawas pemilihan, dan para pemilih. Tujuannya adalah untuk memastikan bahwa seluruh proses pemilihan dapat berjalan dengan lancar, transparan, dan sesuai dengan prosedur yang ditetapkan.

% -----------------------------------------------------------------------------------
% Section: Metode-Metode Pelaksanaan Pemilihan OSIS
% -----------------------------------------------------------------------------------
\section{Metode-Metode Pelaksanaan Pemilihan OSIS}
\label{sec:metode_pemilihan}

Seiring perkembangan zaman, pelaksanaan pemilihan Ketua dan Wakil Ketua OSIS telah berevolusi. Pemilihan metode yang tepat sering kali bergantung pada skala pemilihan, sumber daya, dan infrastruktur yang dimiliki oleh sekolah. Secara umum, terdapat dua pendekatan utama: konvensional dan modern.

\subsection{Metode Konvensional (Manual)}
Metode ini adalah pendekatan tradisional yang mengandalkan proses fisik sepenuhnya. Pemilih menggunakan surat suara kertas untuk mencoblos atau menandai pasangan calon pilihan mereka, yang kemudian dimasukkan ke dalam kotak suara. Setelah periode pemungutan suara berakhir, panitia akan membuka kotak suara dan menghitung setiap suara secara manual di hadapan saksi.

\begin{itemize}
	\item \textbf{Kelebihan:} Prosesnya sangat sederhana dan mudah dipahami oleh semua orang. Tidak memerlukan investasi teknologi yang signifikan, sehingga dapat diterapkan di sekolah dengan keterbatasan fasilitas.
	\item \textbf{Kekurangan:} Proses penghitungan suara memakan waktu yang sangat lama, terutama jika jumlah pemilih besar. Metode ini juga sangat rentan terhadap kesalahan manusia (\textit{human error}), baik dalam proses pencoblosan (suara tidak sah) maupun saat rekapitulasi. Selain itu, penggunaan kertas dalam jumlah besar berdampak buruk bagi lingkungan.
\end{itemize}

\subsection{Metode Modern (Elektronik)}
Metode modern memanfaatkan teknologi untuk mengatasi kelemahan dari metode konvensional, terutama dalam hal kecepatan, akurasi, dan efisiensi.

\subsubsection{Metode Semi-Elektronik}
Metode ini merupakan jembatan antara sistem manual dan digital penuh. Pemilih tetap menggunakan surat suara fisik, namun didesain dalam format khusus seperti \textbf{Lembar Jawaban Komputer (LJK)}. Setelah pemungutan suara, seluruh surat suara akan dipindai (\textit{scanned}) menggunakan mesin pemindai khusus. Perangkat lunak kemudian akan membaca tanda pada LJK dan menghitung hasilnya secara otomatis. Ini secara signifikan mengurangi waktu rekapitulasi dan meminimalisir kesalahan hitung.

\subsubsection{Metode Elektronik Penuh (\textit{E-Voting})}
Metode ini sepenuhnya meninggalkan penggunaan kertas. Seluruh proses, mulai dari pemungutan hingga penghitungan suara, dilakukan secara digital. Pemilih mengakses sebuah aplikasi atau sistem, yang disebut \textbf{Surat Suara Elektronik (SSE)}, melalui perangkat seperti komputer, tablet, atau ponsel pintar. Setelah melakukan autentikasi, pemilih dapat langsung memilih kandidat pada layar. Suara yang masuk akan direkam dan dihitung oleh server secara \textit{real-time}. \textbf{Sistem e-SPPO termasuk dalam kategori ini}, menyediakan solusi \textit{e-voting} yang terintegrasi.

% -----------------------------------------------------------------------------------
% Section: Apa Itu e-SPPO?
% -----------------------------------------------------------------------------------
\section{Apa Itu e-SPPO?}
\label{sec:apa_itu_esppo}

\textbf{Sistem Penyelenggaraan Pemilihan OSIS Elektronik (e-SPPO)} adalah sebuah \textit{software suite} (kumpulan program aplikasi) berbasis web yang komprehensif, dirancang untuk mengelola dan melaksanakan seluruh tahapan pemilihan Ketua dan Wakil Ketua OSIS secara digital.

e-SPPO bukan sekadar aplikasi pemungutan suara, melainkan sebuah ekosistem terpadu yang mencakup manajemen data pemilih, pendaftaran kandidat, proses pemungutan suara yang aman, hingga rekapitulasi dan pelaporan hasil secara otomatis. Sistem ini lahir dari kebutuhan akan sebuah platform pemilihan yang tidak hanya efisien, tetapi juga menjunjung tinggi nilai-nilai transparansi, akuntabilitas, dan keamanan data dalam lingkungan pendidikan.

% -----------------------------------------------------------------------------------
% Section: Fungsi e-SPPO
% -----------------------------------------------------------------------------------
\section{Fungsi e-SPPO}
\label{sec:fungsi_esppo}
Sebagai sebuah sistem terpadu, e-SPPO dirancang untuk menjalankan empat fungsi krusial dalam penyelenggaraan pemilihan OSIS:
\begin{itemize}
	\item \textbf{Pengelolaan Data Terpusat:} e-SPPO menyediakan panel administrasi bagi panitia untuk mengelola seluruh data pemilihan di satu tempat. Ini mencakup impor dan validasi Daftar Pemilih Tetap (DPT), pendaftaran pasangan calon beserta visi, misi, dan foto, serta manajemen akun untuk panitia dan pengawas.
	\item \textbf{Pelaksanaan Pemungutan Suara Elektronik:} Fungsi inti dari e-SPPO adalah menyediakan antarmuka Surat Suara Elektronik (SSE) yang intuitif dan aman bagi pemilih. Sistem memastikan bahwa setiap pemilih yang sah hanya dapat memberikan satu suara, dan kerahasiaan pilihan mereka tetap terjaga.
	\item \textbf{Rekapitulasi Suara Otomatis dan Real-time:} Setiap suara yang masuk melalui SSE akan secara otomatis direkam dan diakumulasikan oleh sistem. Panitia dan pengawas dapat memantau perolehan suara secara langsung (\textit{live count}) melalui dasbor administrasi, menghilangkan kebutuhan akan penghitungan manual yang melelahkan dan berisiko.
	\item \textbf{Pemantauan dan Pelaporan Komprehensif:} Sistem ini memberikan alat bagi panitia untuk memantau tingkat partisipasi pemilih. Setelah pemilihan selesai, e-SPPO dapat secara otomatis menghasilkan berbagai dokumen pelaporan, seperti berita acara hasil pemilihan dan rekapitulasi suara total, yang siap untuk dicetak dan ditandatangani.
\end{itemize}

% -----------------------------------------------------------------------------------
% Section: Sejarah Metode Pemilihan OSIS di SMABAT dan e-SPPO
% -----------------------------------------------------------------------------------
\section{Sejarah Metode Pemilihan OSIS di SMABAT dan e-SPPO}
\label{sec:sejarah_esppo}
e-SPPO adalah puncak dari serangkaian inovasi dan perbaikan sistem pemilihan yang telah diterapkan di SMAN 1 Bati-Bati. Setiap tahap evolusi didasari oleh evaluasi terhadap kekurangan sistem sebelumnya dan keinginan untuk menciptakan proses yang lebih baik.
\begin{enumerate}
	\item \textbf{Metode Manual (1992 -- pertengahan 2010-an):} Sejak pendirian sekolah, pemilihan dilakukan secara konvensional. Awalnya, proses ini dilakukan dengan cara panitia membawa kotak suara berkeliling dari kelas ke kelas. Pada pertengahan dekade 2010-an, sistem ini beralih ke model Tempat Pemungutan Suara (TPS), di mana pemilih datang ke lokasi tertentu untuk memberikan suara. Di akhir periode, seluruh surat suara dihitung secara terpusat. Meskipun prosesnya berjalan, tantangan utama dari metode manual ini adalah waktu dan tenaga yang besar untuk rekapitulasi serta potensi suara tidak sah.
	\item \textbf{Elektronik dengan Google Form (2020-2021):} Pandemi COVID-19 memaksa diterapkannya pembelajaran jarak jauh, yang secara langsung berdampak pada metode pemilihan. Google Form digunakan sebagai solusi darurat untuk memungkinkan pemilihan tetap berlangsung secara elektronik. Namun, metode ini memiliki kelemahan signifikan dalam hal keamanan, validasi pemilih (risiko \textit{double voting}), dan kerahasiaan suara.
	\item \textbf{Aplikasi E-Pilketos (2022-2023):} Seiring implementasi Kurikulum Merdeka, semangat Proyek Penguatan Profil Pelajar Pancasila (P5) dengan tema "Suara Demokrasi" mendorong adopsi teknologi \textit{e-voting}. Aplikasi \textit{open-source} E-Pilketos pertama kali digunakan pada tahun 2022 sebagai tugas proyek P5 oleh kelas X-D dan X-C. Aplikasi ini menjadi demonstrasi teknologi \textit{e-voting} dan pembeda dengan kelas lain yang masih menggunakan metode manual. Atas keberhasilan proyek tersebut, sistem ini kemudian diadopsi secara resmi oleh OSIS dan digunakan untuk pemilihan tahun 2022 dan 2023. Namun, E-Pilketos memiliki beberapa kelemahan fundamental: dikembangkan oleh pihak luar, menggunakan \textit{framework} CodeIgniter versi lama yang tidak kompatibel dengan PHP 8.x, menyimpan data pemilih secara \textit{plain-text} (terbuka), dan menggunakan \textit{packages} yang usang.
	\item \textbf{e-SPPO (2024 - Sekarang):} Untuk mengatasi kelemahan E-Pilketos, e-SPPO dikembangkan secara mandiri (\textit{in-house}) sebagai penggantinya. Versi 1.0.0, yang digunakan untuk pemilihan tahun 2024, membawa perbaikan signifikan: dibangun menggunakan kode PHP native yang kompatibel dengan versi 8.x, memanfaatkan Composer untuk manajemen \textit{packages}, dan mengenkripsi data pemilih dalam suara yang tersimpan menggunakan \texttt{password\_hash()}. Meskipun versi ini masih memiliki banyak kekurangan, hasilnya terbukti jauh lebih memuaskan. Keberhasilan tersebut menjadi fondasi pengembangan versi 2.0.0, yang hadir dengan perbaikan arsitektur, keamanan, dan fitur yang lebih matang.
\end{enumerate}

% -----------------------------------------------------------------------------------
% Section: Kenapa e-SPPO? (Keunggulan)
% -----------------------------------------------------------------------------------
\section{Kenapa e-SPPO? (Keunggulan)}
\label{sec:keunggulan_esppo}
e-SPPO menawarkan serangkaian keunggulan yang menjadikannya solusi ideal untuk pemilihan OSIS di era digital:
\begin{itemize}
	\item \textbf{Keamanan Data Lanjutan:} Keamanan adalah prioritas utama. Pada versi 1.0.0, data pemilih dalam suara dienkripsi menggunakan \texttt{password\_hash()}, yang memberikan peningkatan keamanan secara \textbf{marginal} karena keterbatasan fungsi \textit{hash} Bcrypt untuk tujuan enkripsi. Pada versi 2.0.0, sistem ini beralih ke enkripsi simetris \textbf{AES-256-CBC} yang jauh lebih kuat untuk melindungi data pemilih. Sementara itu, untuk autentikasi akun, sistem tetap mempertahankan mekanisme \textit{hashing} kata sandi yang aman menggunakan \texttt{password\_hash()} dan \texttt{password\_verify()}.
	\item \textbf{Kemudahan Akses dan Penggunaan:} Dibangun dengan \textit{framework} antarmuka modern seperti \textbf{Bootstrap 5} dan \textbf{AdminLTE 3.2}, e-SPPO menjamin pengalaman pengguna yang optimal. Tampilannya responsif, artinya dapat diakses dengan nyaman dari berbagai perangkat, mulai dari PC desktop hingga tablet dan ponsel pintar.
	\item \textbf{Transparansi dan Akuntabilitas:} e-SPPO menjamin bahwa setiap suara terekam secara akurat dan dihitung secara tepat waktu. Selama pemilihan berlangsung, hasil hanya dapat diakses oleh pihak berwenang. Namun, setelah pemilihan selesai, hasil dapat dipublikasikan melalui portal publik untuk transparansi penuh. Sistem ini juga mampu menghasilkan beragam laporan dalam format PDF untuk keperluan dokumentasi dan diseminasi informasi.
	\item \textbf{Fleksibilitas Implementasi:} e-SPPO dirancang untuk sangat adaptif. Sistem ini mendukung pemisahan antara server web dan server basis data, serta dapat diimplementasikan di berbagai lingkungan jaringan: mulai dari intranet yang terisolasi penuh (\textit{air-gapped network}), jaringan lokal dengan akses internet terbatas, hingga di-hosting di \textit{cloud}. Metode pelaksanaannya pun fleksibel, dapat disesuaikan dengan kebutuhan, baik berbasis TPS dengan perangkat khusus (PASSE) di setiap bilik suara maupun diakses secara jarak jauh oleh pemilih.
\end{itemize}

% -----------------------------------------------------------------------------------
% Section: Arsitektur Dasar dan Komponen e-SPPO
% -----------------------------------------------------------------------------------
\section{Arsitektur Dasar dan Komponen e-SPPO}
\label{sec:arsitektur_esppo}
Sebagai aplikasi berbasis web, e-SPPO berjalan di atas arsitektur \textit{client-server}. Pengguna (klien) mengakses aplikasi melalui peramban web, sementara seluruh logika pemrosesan dan penyimpanan data terjadi di sisi server. Infrastruktur server ini dapat diatur dalam dua konfigurasi utama: \textbf{tunggal (\textit{stand-alone})} di mana server web dan basis data berjalan di satu mesin, atau \textbf{terpisah (\textit{split server})} di mana keduanya berjalan di mesin yang berbeda untuk skalabilitas yang lebih baik.

Sistem ini mendukung akses melalui protokol \textbf{HTTP} maupun \textbf{HTTPS}, tergantung pada konfigurasi server web. Penggunaan HTTPS sangat diutamakan, terutama jika sistem dapat diakses dari jaringan internet, untuk menjamin enkripsi data selama transit. Untuk jaringan lokal yang aman dan terisolasi, HTTP dapat menjadi pilihan.

\subsection{Teknologi yang Digunakan}
Fondasi teknologi e-SPPO dipilih untuk memaksimalkan fleksibilitas, keamanan, dan kemudahan pemeliharaan:
\begin{itemize}
	\item \textbf{Back-End:} Menggunakan \textbf{PHP murni} tanpa \textit{framework} tambahan. Pilihan ini memberikan kontrol penuh atas kode, performa yang ringan, dan kemudahan untuk di-deploy di hampir semua layanan hosting web standar.
	\item \textbf{Front-End:} Kombinasi \textit{framework} antarmuka yang berbeda digunakan untuk setiap komponen. \textbf{Bootstrap 5} menjadi dasar untuk \textit{Surat Suara Elektronik (SSE)} dan \textit{Portal Publik}, memberikan tampilan yang bersih dan modern. Sementara itu, \textbf{AdminLTE 3.2} digunakan untuk \textit{Pusminlihdu}, menyediakan templat dasbor yang kaya fitur dan ideal untuk administrasi data.
	\item \textbf{Basis Data:} Sistem berkomunikasi menggunakan protokol \textbf{MySQL}, sehingga kompatibel dengan server basis data \textbf{MySQL} maupun \textbf{MariaDB}, dua sistem manajemen basis data relasional yang paling populer dan andal.
\end{itemize}

\subsection{Komponen Utama e-SPPO Versi 2.0.0}
Secara fungsional, ekosistem e-SPPO terbagi menjadi dua komponen aplikasi utama dan satu komponen pendukung yang saling terintegrasi:

\subsubsection{Surat Suara Elektronik (SSE)}
Ini adalah "wajah" dari e-SPPO yang dilihat oleh \textbf{pemilih}. SSE adalah aplikasi web yang menjadi gerbang bagi pemilih untuk menggunakan hak suaranya. Alur penggunaannya dirancang agar lugas dan sederhana: pemilih membuka aplikasi, melakukan \textit{login} dengan akun yang telah diberikan, memilih salah satu kandidat, mengonfirmasi pilihan tersebut untuk dikirim dan diproses oleh sistem, kemudian melakukan \textit{logout} untuk mengakhiri sesi. Antarmuka SSE sengaja dibuat minimalis untuk memastikan semua pemilih dapat menggunakannya dengan mudah tanpa kebingungan.

\subsubsection{Pusat Administrasi Pemilihan Terpadu (Pusminlihdu)}
Pusminlihdu adalah "ruang kontrol" atau pusat komando dari seluruh kegiatan pemilihan. Ini adalah panel administrasi komprehensif yang hanya dapat diakses oleh \textbf{Panitia Seleksi dan Pemilihan (Panselih)} serta \textbf{Pengawas Pemilihan} dengan hak akses berjenjang. Dari sini, panitia dapat mengelola berbagai jenis data dan konfigurasi, antara lain:
\begin{itemize}
	\item \textbf{Data Administratif:} Mengelola akun untuk pengelola (administrator) dan pengawas, serta memasukkan data pejabat kesiswaan yang akan menandatangani dokumen resmi.
	\item \textbf{Detail Kegiatan Pemilihan:} Mengatur parameter krusial seperti nama kegiatan, jadwal (tanggal mulai dan selesai), jumlah TPS, jumlah Perangkat Akses Surat Suara Elektronik (PASSE), tipe kandidat (tunggal atau berpasangan), mode tampilan SSE, hingga status pemilihan (aktif, selesai, dll).
	\item \textbf{Data Konstituensi dan DPT:} Mengelola data pemilih, mulai dari pengelompokan (misalnya per kelas) hingga Daftar Pemilih Tetap (DPT) yang akan digunakan dalam pemilihan.
	\item \textbf{Data Calon:} Memasukkan dan mengelola informasi mengenai pasangan calon, termasuk nama, foto, visi, dan misi.
	\item \textbf{Fungsi Lainnya:} Memantau jalannya pemilihan secara \textit{real-time}, melihat hasil perolehan suara, dan yang terpenting, membuat berbagai laporan dan berkas pemilihan dalam format PDF.
\end{itemize}

\subsubsection{Portal Web Publik e-SPPO}
Ini adalah komponen pendukung yang berfungsi sebagai etalase informasi pemilihan kepada publik atau seluruh warga sekolah. Portal ini dirancang untuk menampilkan hasil akhir rekapitulasi suara secara transparan setelah seluruh proses pemilihan selesai dan diverifikasi oleh panitia. Selain itu, portal ini juga menjadi tempat di mana laporan dan berkas-berkas lain yang dihasilkan oleh Pusminlihdu dapat diunduh oleh umum, sehingga memperkuat akuntabilitas proses pemilihan.

% --- Akhir dari Bab 1: Pendahuluan ---

% ===================================================================================
% BAB 2: PENYIAPAN INFRASTRUKTUR DAN INSTALASI e-SPPO
% Berisi panduan teknis untuk menyiapkan server dan menginstal aplikasi e-SPPO.
% File ini akan di-include oleh main.tex
% ===================================================================================

\part{Penyiapan Infrastruktur dan Instalasi e-SPPO}
\label{part:penyiapan_sistem}

Bab ini menyajikan panduan teknis yang komprehensif untuk menyiapkan lingkungan server dan melakukan instalasi aplikasi e-SPPO. Proses ini ditujukan untuk administrator sistem atau panitia teknis yang bertanggung jawab atas infrastruktur pemilihan. Keberhasilan instalasi sangat bergantung pada pemenuhan setiap persyaratan dan langkah-langkah konfigurasi yang diuraikan di bawah ini secara cermat.

% -----------------------------------------------------------------------------------
% Section: Persyaratan Sistem dan Aplikasi
% -----------------------------------------------------------------------------------
\section{Persyaratan Sistem dan Aplikasi}
\label{sec:persyaratan_sistem}
Sebelum memulai, pastikan server yang akan digunakan telah memenuhi persyaratan perangkat keras dan lunak berikut.

\subsection{Persyaratan Fisik dan Sistem Operasi}
Meskipun e-SPPO dapat berjalan di komputer biasa (PC desktop atau laptop), penggunaan perangkat keras kelas server sangat disarankan untuk stabilitas dan performa optimal, terutama saat melayani banyak pemilih secara bersamaan. Jika beban kerja dapat dikelola dengan baik (misalnya, jumlah pemilih tidak terlalu besar), komputer biasa dapat difungsikan sebagai server.

\begin{itemize}
	\item \textbf{Perangkat Keras (Hardware):}
	\begin{itemize}
		\item \textbf{Prosesor (CPU):} Diperlukan prosesor yang cukup tangguh untuk memproses data. Direkomendasikan minimal setara dengan \textbf{Intel Core i5 atau i7} generasi modern atau yang lebih tinggi.
		\item \textbf{RAM:}
		\begin{itemize}
			\item Minimum: \textbf{4 GB}
			\item Disarankan: \textbf{8 GB -- 16 GB} untuk kelancaran operasional.
			\item Optimal: \textbf{Di atas 16 GB} untuk skala pemilihan yang besar.
		\end{itemize}
		\item \textbf{Penyimpanan:} Ruang penyimpanan harus cukup untuk menampung:
		\begin{itemize}
			\item Sistem Operasi (Windows/Linux).
			\item \textit{Stack} server (Apache, PHP, MariaDB/MySQL).
			\item Aplikasi pendukung (seperti phpMyAdmin).
			\item Aplikasi e-SPPO beserta seluruh dependensi dan berkas yang akan dihasilkannya (laporan, gambar, dll.).
		\end{itemize}
		\item \textbf{Jaringan:} Koneksi jaringan yang stabil dan andal. Penggunaan koneksi kabel (Ethernet) lebih diutamakan daripada nirkabel.
	\end{itemize}
	\item \textbf{Sistem Operasi (OS) yang Didukung:}
	\begin{itemize}
		\item Microsoft Windows (mulai dari Windows 7 hingga Windows 11).
		\item Microsoft Windows Server (mulai dari versi 2012 hingga 2025).
		\item Distribusi Linux (seperti Ubuntu, CentOS, Debian).
		\item macOS.
	\end{itemize}
\end{itemize}

\subsection{Persyaratan Perangkat Lunak Server}
\begin{itemize}
	\item \textbf{Server Web:} Dioptimalkan untuk \textbf{Apache HTTP Server versi 2.4}.
	\item \textbf{PHP:} Versi minimum yang diperlukan adalah \textbf{PHP 8.2.x}, namun aplikasi ini dirancang dan diuji secara ekstensif pada \textbf{PHP 8.3.x}.
	\item \textbf{Ekstensi PHP Wajib:} Sejumlah ekstensi PHP harus diaktifkan agar semua fitur e-SPPO dapat berjalan. Bagi pengguna yang menginstal PHP secara manual (bukan melalui paket seperti XAMPP), perhatian ekstra pada bagian ini sangat penting. Berikut adalah daftar ekstensi krusial dan fungsinya:
	\begin{description}
		\item[\texttt{mysqli} dan \texttt{pdo\_mysql}] Penting untuk koneksi antara aplikasi PHP dan server basis data MySQL/MariaDB.
		\item[\texttt{curl}] Dibutuhkan untuk komunikasi dengan layanan eksternal atau API di masa depan.
		\item[\texttt{fileinfo}] Digunakan untuk mendeteksi tipe file (MIME type) dari berkas yang diunggah, seperti logo sekolah atau foto kandidat, untuk validasi keamanan.
		\item[\texttt{gd}] Diperlukan untuk pemrosesan gambar, seperti membuat \textit{thumbnail} atau memanipulasi gambar yang diunggah.
		\item[\texttt{intl}] (Internationalization) Penting untuk penanganan format tanggal, waktu, dan angka yang sesuai dengan standar regional (lokal).
		\item[\texttt{mbstring}] (Multi-byte String) Krusial untuk penanganan karakter non-Latin (Unicode) dengan benar, memastikan nama atau data lain dapat ditampilkan tanpa masalah.
		\item[\texttt{exif}] Digunakan untuk membaca metadata dari file gambar, yang mungkin diperlukan untuk validasi tambahan.
		\item[\texttt{openssl}] Diperlukan untuk fungsi-fungsi kriptografi, termasuk enkripsi AES-256-CBC yang digunakan e-SPPO.
		\item[\texttt{yaml}] \textbf{Sangat Penting.} Digunakan untuk membaca semua file konfigurasi e-SPPO (\texttt{*.yml}). Ekstensi ini \textbf{tidak tersedia secara default} dan harus diinstal manual dari PECL:
		\begin{itemize}
			\item Unduh dari: \url{https://pecl.php.net/package/yaml}
			\item Khusus Windows: \url{https://pecl.php.net/package/yaml/2.2.5/windows}
		\end{itemize}
	\end{description}
	\item \textbf{Manajer Dependensi:} \textbf{Composer} harus sudah terinstal secara global di sistem Anda untuk mengelola pustaka-pustaka PHP.
	\item \textbf{Sistem Manajemen Basis Data (DBMS):} Diperlukan server basis data yang mendukung \textbf{protokol MySQL}. Oleh karena itu, Anda dapat menggunakan server \textbf{MySQL} maupun \textbf{MariaDB}.
\end{itemize}

% -----------------------------------------------------------------------------------
% Section: Penyiapan Sistem dan Aplikasi
% -----------------------------------------------------------------------------------
\section{Penyiapan Sistem dan Aplikasi}
\label{sec:penyiapan_aplikasi}
Tahap ini mencakup instalasi dan konfigurasi \textit{stack} server (Apache, PHP, MariaDB/MySQL). Anda dapat menginstalnya satu per satu atau menggunakan paket terintegrasi seperti XAMPP untuk kemudahan.

\subsection{Konfigurasi Wajib Server}
\importantbox{Langkah ini sangat penting. Melewatkan atau salah mengonfigurasi bagian ini dapat menyebabkan aplikasi tidak berjalan atau mengalami galat (\textit{error}) saat digunakan, terutama pada fitur-fitur yang membutuhkan sumber daya lebih.}

Agar e-SPPO berjalan optimal, beberapa pengaturan pada file konfigurasi Apache (\texttt{httpd.conf}) dan PHP (\texttt{php.ini}) \textbf{wajib} disesuaikan.

\subsubsection{Konfigurasi PHP (\texttt{php.ini})}
Beberapa parameter PHP perlu ditingkatkan dari nilai defaultnya untuk mengakomodasi kebutuhan e-SPPO:
\begin{itemize}
	\item \texttt{max\_execution\_time = 210} \\ Waktu eksekusi skrip ditingkatkan untuk mencegah \textit{timeout} saat memproses data dalam jumlah besar, seperti impor DPT atau pembuatan laporan.
	\item \texttt{memory\_limit = 512M} \\ Alokasi memori ditingkatkan agar PHP dapat menangani operasi yang intensif, seperti membaca file Excel yang besar.
	\item \texttt{upload\_max\_filesize = 512M} \\ Batas ukuran file yang diunggah harus cukup besar untuk menampung file DPT atau gambar kandidat.
	\item \texttt{post\_max\_size = 0} \\ (Nilai 0 berarti tidak terbatas) Ukuran data POST ditingkatkan untuk mendukung pengiriman formulir dengan banyak data.
	\item \texttt{date.timezone = Asia/Makassar} \\ Atur zona waktu sesuai lokasi Anda. Ini krusial agar stempel waktu (\textit{timestamp}) pada data suara dan log tercatat dengan benar.
	\item \textbf{Aktifkan Ekstensi:} Hapus tanda titik koma (;) di awal baris ekstensi yang dibutuhkan seperti \texttt{extension=curl}, \texttt{extension=intl}, \texttt{extension=mysqli}, \texttt{extension=openssl}, dan terutama \texttt{extension=yaml}.
\end{itemize}

\subsubsection{Konfigurasi Apache (\texttt{httpd.conf})}
Pastikan Apache dimuat dengan modul yang benar dan dikonfigurasi untuk menangani PHP.
\begin{itemize}
	\item \textbf{Aktifkan Modul:} Pastikan modul seperti \texttt{mod\_rewrite} dan \texttt{mod\_ssl} diaktifkan dengan menghapus tanda pagar (\#) di awal barisnya.
	\item \textbf{Izinkan \texttt{.htaccess}:} Ubah direktif \texttt{AllowOverride} dari \texttt{None} menjadi \texttt{All} untuk direktori root web Anda. Ini penting agar file \texttt{.htaccess} e-SPPO dapat berfungsi.
\end{itemize}
\notebox{Contoh lengkap perubahan yang direkomendasikan pada kedua file tersebut dapat dilihat pada berkas \texttt{httpd.diff.txt} dan \texttt{php.diff.txt} yang disertakan dalam lampiran panduan ini.}

% -----------------------------------------------------------------------------------
% Section: Instalasi e-SPPO
% -----------------------------------------------------------------------------------
\section{Instalasi e-SPPO}
\label{sec:instalasi_esppo}
\subsection{Memperoleh dan Memasang Kode Sumber}
\begin{enumerate}
	\item Unduh rilis terbaru e-SPPO dari repositori resmi: \url{https://github.com/IndoRizkikali/esppo-pilketos}.
	\item Ekstrak file arsip \texttt{.zip} yang telah diunduh.
	\item Salin seluruh isi folder hasil ekstraksi ke dalam direktori root server web Anda (misalnya \texttt{C:/xampp/htdocs/esppo}).
\end{enumerate}

\subsection{Instalasi Dependensi Menggunakan Composer}
\begin{enumerate}
	\item Buka terminal atau \textit{command prompt}.
	\item Arahkan direktori aktif ke folder instalasi e-SPPO \\(misal: \texttt{cd C:/xampp/htdocs/esppo}).
	\item Jalankan perintah: \texttt{composer install}. Tunggu hingga proses selesai.
\end{enumerate}

% -----------------------------------------------------------------------------------
% Section: Pengaturan Awal e-SPPO
% -----------------------------------------------------------------------------------
\section{Pengaturan Awal e-SPPO}
\label{sec:pengaturan_awal}
\warningbox{Semua file konfigurasi di dalam direktori \texttt{confs/} harus dibuat dari file contoh (\texttt{*.example.yml}) dan diisi dengan lengkap. Setiap variabel di dalamnya digunakan oleh aplikasi, dan jika ada yang kosong atau salah format, dapat menyebabkan aplikasi gagal berfungsi.}

\subsection{Pembuatan dan Konfigurasi Basis Data}
\begin{enumerate}
	\item \textbf{Buat Basis Data:} Menggunakan phpMyAdmin, buat basis data baru (misal: \texttt{sppoe\_2025}).
	\item \textbf{Buat Pengguna (Direkomendasikan):} Buat akun pengguna basis data baru yang hanya memiliki hak akses penuh ke basis data tersebut.
	\item \textbf{Impor Struktur Tabel:} Impor file \texttt{sppoe\_db\_tables.sql} ke dalam basis data yang baru dibuat.
	\item \textbf{Konfigurasi Koneksi:} Salin \texttt{confs/db\_config.example.yml} menjadi \texttt{db\_config.yml}, lalu edit isinya dengan informasi koneksi basis data Anda.
\end{enumerate}

\subsection{Pengaturan Infrastruktur Kriptografi}
\begin{enumerate}
	\item \textbf{Buat Kunci Enkripsi:} Buka perkakas web \texttt{tools/esppo-keymaker.php} melalui peramban. Salin kunci yang dihasilkan.
	\item \textbf{Konfigurasi Kunci:} Salin \texttt{confs/crypto\_config.example.yml} menjadi \texttt{crypto\_config.yml}, lalu tempelkan kunci yang telah Anda buat.
\end{enumerate}

\subsection{Pengaturan Data Sekolah}
\begin{enumerate}
	\item Salin file \texttt{confs/school\_config.example.yml} menjadi \texttt{school\_config.yml}.
	\item Buka file \texttt{school\_config.yml} dan \textbf{isi semua informasi yang diminta secara lengkap dan akurat}.
\end{enumerate}
\begin{verbatim}
	# Contoh lengkap file school_config.yml yang harus diisi
	school_config:
	- nama_sekolah: "SMA Negeri 1 Bati-Bati"
	- npsn_sekolah: 30300701
	- alamat_sekolah:
	- alamat: "Jl. A. Yani, KM. 33.3"
	- kecamatan: "Bati-Bati"
	- dati_2: "Kabupaten Tanah Laut"
	- dati_1: "Kalimantan Selatan"
	- kode_pos: 70852
	- no_telepon: 05114779299
	- email_sekolah: "sman1bati.bati@gmail.com"
	- website_sekolah: "https://www.sman1batibati.sch.id"
	- kepala_sekolah:
	- nama: "Taslim, S.Pd., M.Pd."
	- nip_kepsek: 197106021998021005
	- tahun_ajaran:
	- tahun_ajaran: "2025/2026"
	- semester: "Ganjil"
	- logo:
	- logo_sekolah: "Logo_SMABAT_Besar.png"
	- logo_osis: "Logo_OSIS_SMABAT.png"
	- url_akses_esppo:
	- url_sse:
	- "https://192.168.0.11/esppo/"
	- url_pspo:
	- "https://192.168.0.11/esppo/"
\end{verbatim}

\notebox{Setelah semua langkah di atas selesai dan semua file konfigurasi telah diisi dengan benar, aplikasi e-SPPO Anda siap untuk diakses. Langkah selanjutnya adalah melakukan login awal ke Pusminlihdu untuk memulai pengelolaan data pemilihan.}

% --- Akhir dari Bab 2: Penyiapan Infrastruktur dan Instalasi e-SPPO ---

\part{Panduan Panitia -- Pusat Administrasi Pemilihan Terpadu (PUSMINLIHDU)}
% ===================================================================================
% BAB 4: PANDUAN PENGGUNAAN e-SPPO UNTUK PEMILIH -- SURAT SUARA ELEKTRONIK (SSE)
% Berisi panduan lengkap bagi pemilih untuk menggunakan aplikasi SSE.
% File ini akan di-include oleh main.tex
% ===================================================================================

\part{Panduan Penggunaan e-SPPO untuk Pemilih -- Surat Suara Elektronik (SSE)}
\label{part:penggunaan_sse}
Bab ini ditujukan khusus bagi Anda, para pemilih. Di sini, Anda akan menemukan panduan lengkap tentang cara menggunakan aplikasi Surat Suara Elektronik (SSE) untuk memberikan hak suara Anda secara digital. Prosesnya dirancang agar sederhana, cepat, dan aman.

% -----------------------------------------------------------------------------------
% Section: Pengenalan Surat Suara Elektronik (SSE)
% -----------------------------------------------------------------------------------
\section{Pengenalan Surat Suara Elektronik (SSE)}
\label{sec:pengenalan_sse}
Surat Suara Elektronik (SSE) adalah antarmuka berbasis web yang berfungsi sebagai "kertas suara digital" Anda dalam sistem e-SPPO. Melalui SSE, Anda dapat melihat daftar kandidat, mempelajari visi dan misi mereka, serta memberikan pilihan Anda dengan beberapa klik saja.

\subsection{Keunggulan SSE bagi Pemilih}
\begin{itemize}
	\item \textbf{Mudah Diakses dan Digunakan:} Dengan antarmuka yang modern dan responsif, SSE dapat diakses dari berbagai perangkat (komputer, laptop, tablet, ponsel) dan sangat mudah digunakan.
	\item \textbf{Proses Cepat:} Alur penggunaan yang sederhana memungkinkan Anda memberikan suara hanya dalam hitungan menit, mengurangi waktu antrean di TPS.
	\item \textbf{Kerahasiaan Terjamin:} Pilihan Anda direkam secara anonim. Identitas Anda dienkripsi dengan teknologi \textbf{AES-256-CBC}, sehingga tidak dapat dihubungkan langsung dengan suara yang Anda berikan.
	\item \textbf{Akurasi Terjamin:} Sistem menghitung suara secara otomatis, menghilangkan risiko kesalahan hitung yang biasa terjadi pada metode manual.
\end{itemize}

% -----------------------------------------------------------------------------------
% Section: Cara Kerja dan Keamanan SSE
% -----------------------------------------------------------------------------------
\section{Cara Kerja dan Keamanan SSE}
\label{sec:cara_kerja_sse}
Memahami cara kerja SSE dapat memberikan Anda keyakinan bahwa suara Anda aman dan tercatat dengan benar.

\subsection{Alur Proses Teknis}
Saat Anda menggunakan SSE, berikut adalah proses teknis yang terjadi di belakang layar, sesuai dengan kode programnya:
\begin{enumerate}
	\item \textbf{Validasi Kredensial (\texttt{login.php}):} Ketika Anda memasukkan Nama Akun dan Kata Sandi, sistem akan mencocokkan \textit{hash} kata sandi yang Anda masukkan dengan yang tersimpan di basis data menggunakan fungsi \texttt{password\_verify()}. Sistem juga akan memeriksa apakah status pemilihan sedang berlangsung dan apakah akun Anda belum digunakan untuk memilih.
	\item \textbf{Pencatatan Kehadiran:} Jika login berhasil, sistem akan langsung menyimpan data kehadiran Anda (ID unik pemilih, jenis kelamin, dan konstituensi) ke dalam tabel \texttt{data\_kehadiran}. Proses ini terjadi \textbf{sebelum} Anda memberikan suara, sehingga memisahkan data kehadiran dari data pilihan.
	\item \textbf{Penyimpanan Sesi:} Informasi penting seperti ID unik, nama, jenis kelamin, dan \textit{Initialization Vector} (IV) untuk enkripsi disimpan sementara dalam sesi (\texttt{\$\_SESSION}) untuk digunakan pada proses selanjutnya.
	\item \textbf{Proses Perekaman Suara (\texttt{ballot\_process.php}):} Ini adalah bagian paling krusial.
	\begin{itemize}
		\item \textbf{Transaksi Atomik:} Seluruh proses penyimpanan suara dijalankan dalam sebuah transaksi basis data. Ini berarti semua langkah harus berhasil; jika satu langkah gagal, seluruh proses akan dibatalkan (\textit{rollback}) untuk mencegah data yang tidak konsisten.
		\item \textbf{Penguncian Data:} Sistem mengunci baris data Anda di tabel pemilih (\texttt{FOR UPDATE}) untuk mencegah upaya pemilihan ganda secara bersamaan (\textit{race condition}).
		\item \textbf{Enkripsi Identitas:} ID unik Anda dienkripsi menggunakan kunci rahasia dan IV unik milik akun Anda dengan algoritma \textbf{AES-256-CBC}. Hasil enkripsi inilah yang disimpan bersama suara Anda, bukan ID asli.
		\item \textbf{Pembuatan ID Suara Unik:} Sistem menghasilkan sebuah ID unik untuk surat suara Anda yang digabungkan dari ID terenkripsi Anda dan ID kandidat yang dipilih.
		\item \textbf{Penyimpanan Suara:} Data suara (ID suara, nomor urut kandidat, ID pemilih terenkripsi, dll.) disimpan ke tabel \texttt{data\_suara}.
		\item \textbf{Pembaruan Status:} Status Anda di tabel \texttt{data\_pemilih} diubah menjadi `SUDAH\_MEMILIH`.
		\item \textbf{Commit Transaksi:} Jika semua langkah sukses, perubahan disimpan secara permanen.
	\end{itemize}
	\item \textbf{Penghancuran Sesi:} Setelah suara berhasil direkam, sesi Anda akan langsung dihancurkan (\texttt{session\_destroy()}) untuk keamanan. Anda akan otomatis keluar dari sistem.
\end{enumerate}

% -----------------------------------------------------------------------------------
% Section: Ketentuan Penggunaan SSE
% -----------------------------------------------------------------------------------
\section{Ketentuan Penggunaan SSE}
\label{sec:ketentuan_sse}
\begin{itemize}
	\item \textbf{Satu Akun, Satu Suara:} Setiap akun pemilih hanya dapat digunakan untuk memberikan suara satu kali.
	\item \textbf{Jaga Kerahasiaan Akun:} Jangan pernah membagikan Nama Akun dan Kata Sandi Anda kepada siapa pun.
	\item \textbf{Gunakan Perangkat yang Aman:} Jika memilih dari TPS, gunakan hanya Perangkat Akses Surat Suara Elektronik (PASSE) yang disediakan. Jangan mengutak-atik perangkat tersebut.
	\item \textbf{Laporkan Kendala:} Jika terjadi galat atau masalah teknis, segera laporkan kepada panitia pemilihan yang bertugas.
\end{itemize}

% -----------------------------------------------------------------------------------
% Section: Panduan Langkah Demi Langkah
% -----------------------------------------------------------------------------------
\section{Panduan Langkah Demi Langkah}
\label{sec:langkah_penggunaan_sse}
Ikuti langkah-langkah berikut untuk memberikan suara Anda melalui SSE.

\subsection{Langkah 1: Mengakses Halaman Login}
Buka peramban web pada perangkat Anda, lalu masukkan alamat URL SSE yang diberikan oleh panitia. Anda akan disambut oleh halaman login yang menampilkan nama kegiatan pemilihan yang sedang berlangsung.
\begin{figure}[h]
	\centering
	\includegraphics[width=1\textwidth]{SSE_1.png}
	\caption{Halaman Login Aplikasi SSE.}
	\label{fig:sse_login}
\end{figure}

\subsection{Langkah 2: Masuk (Login) ke Akun Anda}
Masukkan \textbf{Nama Akun Pemilih} dan \textbf{Kata Sandi} Anda pada kolom yang tersedia, lalu klik tombol \textbf{Masuk}.
\begin{figure}[h]
	\centering
	\includegraphics[width=1\textwidth]{SSE_2.png}
	\caption{Mengisi Kredensial Akun Pemilih.}
	\label{fig:sse_login_filled}
\end{figure}

\subsection{Langkah 3: Memilih Kandidat}
Setelah berhasil login, Anda akan masuk ke halaman surat suara utama.
\begin{enumerate}
	\item \textbf{Pelajari Kandidat:} Luangkan waktu untuk melihat foto, nama, serta visi dan misi setiap kandidat yang ditampilkan.
	\item \textbf{Tentukan Pilihan:} Klik pada tombol radio (lingkaran) atau pada area kartu kandidat pilihan Anda. Tanda centang biru akan muncul untuk menandakan pilihan Anda.
\end{enumerate}
\begin{figure}[h]
	\centering
	\includegraphics[width=1\textwidth]{SSE_4.png}
	\caption{Memilih Salah Satu Kandidat.}
	\label{fig:sse_selecting_candidate}
\end{figure}

\subsection{Langkah 4: Mengirim dan Mengonfirmasi Pilihan}
\begin{enumerate}
	\item Setelah yakin, klik tombol biru \textbf{Kirim Suara Saya}.
	\item Sebuah kotak dialog konfirmasi akan muncul. Ini adalah kesempatan terakhir Anda untuk memeriksa kembali pilihan.
	\item Klik \textbf{Ya - Kirim Suara} untuk melanjutkan, atau \textbf{Tidak - Kembali ke Surat Suara} untuk batal.
\end{enumerate}
\begin{figure}[h]
	\centering
	\includegraphics[width=1\textwidth]{SSE_5.png}
	\caption{Kotak Dialog Konfirmasi Pilihan.}
	\label{fig:sse_confirm_vote}
\end{figure}
\importantbox{Setelah Anda menekan tombol "Ya - Kirim Suara", pilihan Anda bersifat \textbf{final} dan tidak dapat diubah kembali.}

\subsection{Langkah 5: Selesai dan Keluar (Logout)}
Jika suara berhasil dikirim, Anda akan melihat halaman konfirmasi akhir.
\begin{figure}[h]
	\centering
	\includegraphics[width=1\textwidth]{SSE_6.png}
	\caption{Halaman Konfirmasi Akhir.}
	\label{fig:sse_vote_recorded}
\end{figure}
Klik tombol \textbf{Selesai \& Kembali ke Login} untuk keluar sepenuhnya dari akun Anda.
\warningbox{Langkah ini \textbf{wajib} dilakukan, terutama jika Anda menggunakan perangkat umum di TPS, untuk melindungi privasi dan mencegah penyalahgunaan akun.}

% -----------------------------------------------------------------------------------
% Section: Rekomendasi dan Informasi Tambahan
% -----------------------------------------------------------------------------------
\section{Rekomendasi untuk Pengguna SSE}
\label{sec:rekomendasi_sse}
\begin{itemize}
	\item \textbf{Gunakan Koneksi Stabil:} Pastikan Anda terhubung ke jaringan internet atau LAN yang stabil saat memilih untuk menghindari kegagalan pengiriman suara.
	\item \textbf{Jangan Terburu-buru:} Baca semua informasi dengan teliti sebelum membuat pilihan dan mengirimkan suara.
	\item \textbf{Jika Ragu, Bertanya:} Jangan ragu untuk bertanya kepada panitia jika ada hal yang tidak Anda mengerti selama proses pemilihan.
\end{itemize}

\section{Informasi Lainnya}
\label{sec:info_lainnya_sse}
\begin{itemize}
	\item \textbf{Pesan Galat (\textit{Error Messages}):} Jika Anda menemui pesan galat seperti "Nama Akun atau Kata Sandi salah" atau "Akun sudah digunakan", ikuti petunjuk yang diberikan atau hubungi panitia.
	\item \textbf{Reset Akun:} Dalam kasus yang sangat jarang terjadi, seperti salah memberikan akun kepada pemilih atau terjadi kesalahan teknis, panitia memiliki wewenang untuk mereset status pemilihan Anda setelah melalui verifikasi yang ketat.
\end{itemize}

\notebox{Selamat! Anda telah berhasil menggunakan hak suara Anda secara elektronik melalui e-SPPO. Terima kasih atas partisipasi Anda dalam menyukseskan pemilihan yang jujur, adil, dan modern.}

% --- Akhir dari Bab 4: Panduan Penggunaan SSE ---

\part{Pengaturan dan Layanan Lanjutan e-SPPO}
\include{./TeX_files/06-pemecahanmasalah}
\part{Informasi Tambahan dan Serba-Serbi}
\include{./TeX_files/08-penutupdanbantuan}

\part{Apendiks A -- Format Berkas Konfigurasi YAML}
\section{school\_config.yml}
\end{document}